\chapter{Implementasi Algoritma}

Dalam permainan Queens LinkedIn, tujuannya adalah menempatkan satu \textit{queen} pada setiap baris di papan berukuran $N \times M$ yang terbagi menjadi daerah warna, dengan syarat: tepat satu queen pada setiap baris, kolom, dan daerah warna, serta tidak boleh ada dua queen yang saling bersebelahan (termasuk diagonal). Program ini ditulis dalam bahasa C++ dengan GUI menggunakan library SFML 3.0 dan mengimplementasikan algoritma Brute Force.

\section{Algoritma Brute Force}

Algoritma \textit{brute force} yang diimplementasikan bersifat murni, yaitu mengenumerasi seluruh kemungkinan penempatan queen tanpa heuristik apapun.

Berikut adalah langkah-langkah detail dari algoritma brute force yang diimplementasikan pada fungsi \texttt{solve\_bf()}:

\begin{enumerate}
    \item Program menerima masukan berupa \textit{grid} berukuran $N \times M$ yang dibaca dari file \texttt{.txt}. Setiap huruf pada grid merepresentasikan daerah warna.
    \item Dibuat sebuah \textit{array} \texttt{curr} berukuran $N$, diinisialisasi dengan $[0, 0, \dots, 0]$. Elemen \texttt{curr[i]} menyatakan kolom tempat queen diletakkan pada baris ke-$i$.
    \item Untuk setiap konfigurasi \texttt{curr}, dilakukan validasi terhadap seluruh syarat permainan:
    \begin{itemize}
        \item \textbf{Syarat kolom}: Tidak boleh ada dua queen pada kolom yang sama, diperiksa dengan membandingkan \texttt{curr[i]} terhadap \texttt{curr[prev]} untuk seluruh baris sebelumnya.
        \item \textbf{Syarat warna}: Tidak boleh ada dua queen pada daerah warna yang sama, diperiksa dengan membandingkan \texttt{grid[i][curr[i]]} terhadap \texttt{grid[prev][curr[prev]]}.
        \item \textbf{Syarat adjacency}: Queen pada baris berurutan tidak boleh berada pada kolom yang bersebelahan, diperiksa dengan kondisi $|curr[i] - curr[i-1]| \leq 1$.
    \end{itemize}
    \item Jika seluruh syarat terpenuhi, konfigurasi disimpan sebagai solusi.
    \item Jika tidak valid, konfigurasi di-\textit{increment}: \texttt{curr[N-1]} dinaikkan 1; jika melebihi $M$, di-\textit{reset} ke 0 dan \texttt{curr[N-2]} dinaikkan, dan seterusnya.
    \item Proses berulang hingga solusi ditemukan atau seluruh $M^N$ konfigurasi habis diperiksa.
\end{enumerate}

Algoritma ini memiliki kompleksitas waktu \textbf{$O(M^N \cdot N)$} karena setiap konfigurasi memerlukan $O(N)$ untuk validasi.

\begin{algorithm}[H]
\caption{Algoritma Brute Force Queens}
\begin{algorithmic}[1]
\Procedure{SolveBruteForce}{$grid, N, M$}
    \State $curr \gets [0, 0, \dots, 0]$ \Comment{Array berukuran $N$}
    \While{\textbf{true}}
        \State $iter \gets iter + 1$
        \State $valid \gets$ \textbf{true}
        \For{$i \gets 0$ \textbf{to} $N-1$}
            \State $c \gets curr[i]$
            \If{$c \geq M$} $valid \gets$ \textbf{false}; \textbf{break} \EndIf
            \For{$prev \gets 0$ \textbf{to} $i-1$}
                \If{$curr[prev] = c$} $valid \gets$ \textbf{false}; \textbf{break} \EndIf \Comment{Kolom sama}
                \If{$grid[prev][curr[prev]] = grid[i][c]$} $valid \gets$ \textbf{false}; \textbf{break} \EndIf \Comment{Warna sama}
            \EndFor
            \If{$i > 0$ \textbf{and} $|c - curr[i-1]| \leq 1$} $valid \gets$ \textbf{false}; \textbf{break} \EndIf \Comment{Adjacency}
        \EndFor
        \If{$valid$} \Return $curr$ \Comment{Solusi ditemukan}
        \EndIf
        \State \Comment{Increment ke konfigurasi berikutnya}
        \State $idx \gets N - 1$
        \While{$idx \geq 0$}
            \State $curr[idx] \gets curr[idx] + 1$
            \If{$curr[idx] < M$} \textbf{break} \EndIf
            \State $curr[idx] \gets 0$; $idx \gets idx - 1$
        \EndWhile
        \If{$idx < 0$} \Return \texttt{None} \Comment{Tidak ada solusi}
        \EndIf
    \EndWhile
\EndProcedure
\end{algorithmic}
\end{algorithm}

\section{Fitur Tambahan Program}

Selain algoritma pencarian, program juga mengimplementasikan:

\begin{enumerate}
    \item \textbf{GUI (SFML 3.0)}: Panel kontrol dengan tombol-tombol interaktif (BROWSE, LOAD, SOLVE, STOP, SAVE TXT, SAVE IMG), pemilihan algoritma, dan slider kecepatan visualisasi.
    \item \textbf{Live Visualization}: Setiap iterasi pencarian divisualisasikan secara \textit{real-time} pada papan GUI. Kecepatan dapat diatur melalui slider (TURBO / FAST / SLOW).
    \item \textbf{Validasi Input}: Program memverifikasi:
    \begin{itemize}
        \item Seluruh baris memiliki panjang yang sama
        \item Karakter hanya boleh berupa huruf A--Z
        \item Jumlah region warna unik harus sama dengan jumlah baris ($N$)
    \end{itemize}
    \item \textbf{Output as Text}: Solusi dapat disimpan sebagai file \texttt{.txt} melalui tombol SAVE TXT, lengkap dengan waktu pencarian dan jumlah iterasi.
    \item \textbf{Output as Image}: Solusi dapat disimpan sebagai file PNG melalui tombol SAVE IMG.
\end{enumerate}
